% =========================================================================
% limitations.tex
% =========================================================================

\section{Limitations}
\label{sec:limitations}

I state the limitations of this study explicitly.

\paragraph{Head dimension mismatch.}
All experiments use GPT-2 Medium with $\dhead = 64$.
The TurboQuant~\cite{zandieh2025turboquant} and
PolarQuant~\cite{han2025polarquant} papers target Llama-3.1-8B,
Mistral-7B, and Qwen-2.5-7B, which have $\dhead = 128$.
Absolute perplexity numbers are not comparable across these settings.
The relative rankings among schemes within this $\dhead = 64$ evaluation
are valid, but whether the Adaptive QJL correction improves or is
unnecessary at $\dhead = 128$ remains untested.
Cross-dimension generalization of the Adaptive QJL correction to
$\dhead = 128$ was not evaluated in this study and is reserved for
future work.

\paragraph{No fused CUDA kernels.}
All quantization is implemented in pure PyTorch.
The WHT has theoretical complexity $O(d \log d)$ versus $O(d^2)$ for
QR rotation, but in practice the unfused PyTorch WHT is slower than
naive quantization due to the overhead of sequential butterfly
operations.
% Source: results/rotation_profile.json; README.md → Limitation bullet 2
This finding is documented in the rotation profiling results
(\texttt{results/rotation\_profile.json}) and means that the latency
advantages claimed in TurboQuant cannot be realized without kernel
fusion.

\paragraph{Training-scale mismatch for MC-RNN.}
The from-scratch MC-RNN training experiment uses a 15M-parameter
linear-attention model trained on WikiText-103.
The Memory Caching paper~\cite{behrouz2026mc} trains models of
760M to 1.3B parameters on 30--100 billion tokens.
The experiment reproduces the core finding (GRM perplexity 21.9 versus
standard linear attention perplexity 81.9) but at a scale that is
orders of magnitude smaller than the original.
% Source: results/mc_training_results.json →
%   perplexity: grm = 21.904, standard = 81.946;
%   README.md line 236 (scale comparison)

\paragraph{Context window proxy limit.}
A central limitation of this study is the mismatch between the intended application of KV cache quantization and the evaluation model. Quantization is designed for long-context retrieval spanning tens or hundreds of thousands of tokens. However, the evaluation relies on GPT-2 Medium, which is hard-capped at a context window of 1{,}024 tokens. This study functions as a proxy evaluation, focusing on the geometric and statistical failure modes of quantization in low-dimensional spaces ($\dhead = 64$) rather than validating long-context engineering stability at scale.

\paragraph{LongBench-lite truncation.}
LongBench-lite scores are low across all schemes (passage QA
$\leq 25\%$, summarization ROUGE-1 $\leq 12\%$, code completion
$\leq 5\%$) because GPT-2 Medium's 1{,}024-token context limit
truncates most examples.
% Source: results/longbench_lite.json →
%   Baseline: passage_qa=25.0, summarization=11.99, code=0.0
These scores reflect the model's context limitation, not the quality
of any quantization scheme.

\paragraph{Single-model evaluation.}
I evaluate on a single model (GPT-2 Medium).
Conclusions about scheme rankings may not transfer to architectures with
different attention patterns, embedding distributions, or training data.
